% ============================================================================
% 第 1 章 - 项目介绍
% ============================================================================
\chapter{项目介绍}

这是一个从底层向上构建的学习项目。我们不依赖高级语言编译器和操作系统提供的便利，而是亲自去触碰：
\begin{itemize}
    \item \textbf{芯片}（Chip）：用 Verilog 写电路，从简单的字节加法器，到 UART 串口收发器，再到自制 CPU
    \item \textbf{系统}（ISA + QEMU）：用 AArch64 汇编直接写用户态程序，然后把代码加载到 QEMU 仿真的裸机环境里运行
    \item \textbf{算法}：把学到的底层能力，用 C 实现标准算法和数据结构
    \item \textbf{AI}：用 CUDA 并行计算，体会 GPU 与 CPU 协作的底层原理
\end{itemize}

整个项目的核心哲学是：\textbf{自己动手做一遍，才真正理解了这一层在做什么。}

\section{目录结构}

项目顶层目录组织如下：

\vspace{0.5cm}

\noindent\textbf{ISA/} \quad 用户态汇编程序区。
它包含一个完整的 AArch64 汇编程序：手工实现了堆内存分配器、字符串管理、
命令解析器。所有 I/O 都是通过 ARM 的 \texttt{svc} 指令直接触发系统调用。
主要文件：
\begin{itemize}
    \item \texttt{hello.s} — 主程序入口，循环读取、复制、存储、打印用户输入
    \item \texttt{src/exit.s} — \texttt{sys\_write}、\texttt{sys\_read}、\texttt{exit} 的系统调用封装
    \item \texttt{src/memory.s} — \texttt{memory\_init}、\texttt{memory\_alloc}，一个 4 KB 堆的分配器
    \item \texttt{src/memoryfree.s} — \texttt{memory\_free}，带相邻块合并
    \item \texttt{src/memorycopy.s} — \texttt{copy\_string}，从输入缓冲区复制到新分配的内存
    \item \texttt{src/string\_array.s} — 平行数组存储字符串指针与长度，按序号存取
    \item \texttt{src/str\_parse\_cmd.s} — \texttt{find\_two\_spaces}，扫描字符串找到前两个空格的位置
    \item \texttt{src/test.s} — 测试函数：内存读写、系统调用的测试
    \item \texttt{src/memory.inc} — 常量定义（如 \texttt{HEAP\_SIZE = 4096}）
    \item \texttt{run.sh} — 一键编译并启动 lldb 调试
\end{itemize}

\vspace{0.5cm}

\noindent\textbf{QEMU/} \quad 裸机 AArch64 程序区。
不使用任何操作系统，直接在 QEMU 仿真的 \texttt{virt} 机器上运行，通过内存映射
I/O 访问 PL011 UART 输出字符串。主要文件：
\begin{itemize}
    \item \texttt{start.s} — 启动代码：设置 UART 基地址，逐字节输出，然后进入 \texttt{wfe} 等待
    \item \texttt{run.sh} — 用 Homebrew 的 clang/LLVM 交叉编译、链接、启动 QEMU
    \item \texttt{bin/kernel.elf} — 编译产物（链接地址 \texttt{0x40000000}）
\end{itemize}

\vspace{0.5cm}

\noindent\textbf{Chip/} \quad Verilog 硬件设计区。
用寄存器传输级（RTL）代码描述硬件电路。主要文件：
\begin{itemize}
    \item \texttt{src/adder.v} — 64 位字节加法器：把 8 个字节逐位相加，验证移位与位提取
    \item \texttt{src/uart\_rx.v} — UART 接收器：5 状态状态机，按波特率采样 RX 引脚
    \item \texttt{src/uart\_tx.v} — UART 发送器：4 状态状态机，逐位把字节送到 TX 引脚
    \item \texttt{design.md} — 硬件设计思路文档
\end{itemize}

\section{阅读顺序建议}

如果你从底层开始学习，推荐的阅读顺序是：

\begin{enumerate}
    \item 先读本章的 \textbf{ISA}（第 2 章）：你已经在 macOS 上跑过用户态汇编，
        从系统调用和内存管理入手最自然。
    \item 再读 \textbf{QEMU 裸机}（第 3 章）：体会没有操作系统时，程序是怎么
        与硬件直接通信的。
    \item 最后读 \textbf{Chip 硬件设计}（第 4 章）：当你理解了软件如何使用硬件，
        就可以反过来设计硬件本身了。
\end{enumerate}

如果你对硬件更感兴趣，也可以从第 4 章入手，再回头看软件层是怎么调用它的。
